% 正方形绕一边旋转。左图沿光线看，右图沿旋转轴看。
\begin{tikzpicture}[x=1cm,y=1cm,font=\small,>=Stealth]
  \draw[rounded corners,draw=line,fill=pale!35] (0,0) rectangle (15.7,5.0);
  \draw[line] (7.45,.25)--(7.45,4.75);
  \node[anchor=west,font=\bfseries] at (.32,4.62) {沿着平行光看卡片};
  \node[anchor=west,font=\bfseries] at (7.75,4.62) {从侧面看倾斜的卡片};
  \foreach \cx/\width/\labeltext in {2.05/2/正对光源,5.48/1/转过 $60^\circ$} {
    \draw[line,dashed] (\cx-1,1.42) rectangle (\cx+1,3.42);
    \draw[fill=blue!17,draw=blue,thick] (\cx-\width/2,1.42) rectangle (\cx+\width/2,3.42);
    \draw[<->,blue] (\cx-\width/2,1.14)-- node[below] {$\width$ cm} (\cx+\width/2,1.14);
    \node at (\cx,3.82) {\labeltext};
  }
  \draw[<->,muted] (.76,1.42)--node[left] {$2$} (.76,3.42);
  \node at (2.05,.47) {$4\,\mathrm{cm}^2$};
  \node at (5.48,.47) {$2\,\mathrm{cm}^2$};
  \begin{scope}[shift={(10.7,1.02)}]
    \coordinate (a) at (0,0);
    \coordinate (b) at (1.50,2.598);
    \coordinate (c) at (1.50,0);
    \draw[fill=blue!7,draw=none] (a)--(b)--(c)--cycle;
    \draw[dashed,muted] (b)--(c);
    \draw[thick,blue] (a)--(b);
    \draw[thick,blue] (a)--(c);
    \draw[muted] (1.36,0)--(1.36,.14)--(1.50,.14);
    \draw[blue] (.57,0) arc[start angle=0,end angle=60,radius=.57];
    \node at (.95,.44) {$60^\circ$};
    \node[below] at (.75,-.05) {从光源看到的宽度};
    \node[anchor=west] at (1.66,1.7) {卡边长 $2$ cm};
    \draw[->,thick] (a)--(0,2.35) node[above] {指向光源};
    \draw[->,thick,blue] (a)--(-1.84,1.06) node[above left,inner sep=1pt] {$N$};
    \draw[muted] (0,.83) arc[start angle=90,end angle=150,radius=.83];
    \node at (-.72,1.24) {$\theta$};
  \end{scope}
\end{tikzpicture}
